\subsection{Градиентные методы первого порядка}

\subsubsection*{Откуда берётся антиградиент}

Пусть строится последовательность $\{x_k\}$ и пробная точка
\[
  x = x_k + \alpha p_k .
\]
Разложение по формуле Тейлора:
\[
  f(x) = f(x_k) + \alpha\, \nabla\T f(x_k)\, p_k
       + \tfrac12 \alpha^2\, p_k\T H(\bar x_k)\, p_k ,
  \qquad
  \bar x_k = x_k + \theta\,(x - x_k), \quad \theta \in [0,1] .
\]

При \emph{малом} $\alpha$ знак приращения определяется линейным членом,
поэтому
\[
  f(x) < f(x_k)
  \qquad\Longleftarrow\qquad
  \nabla\T f(x_k)\, p_k < 0 .
\]

\inputfig{figures/fig-descent-cone}

Для градиентных методов первого порядка берут
\[
  \boxed{\ p_k = -\nabla f(x_k)\ } \qquad \text{— антиградиент},
  \qquad \norm{\nabla f(x_k)} \ne 0 .
\]

\begin{idea}
Проверка условия убывания для этого выбора тривиальна:
$\nabla\T f(x_k)\,p_k = -\norm{\nabla f(x_k)}^2 < 0$, пока градиент не нулевой.
Более того, среди всех направлений фиксированной длины антиградиент даёт
наибольшее по модулю отрицательное значение $\nabla\T f(x_k)\,p_k$ — отсюда и
название «наискорейший спуск» для варианта с точным шагом.
\end{idea}

\subsubsection*{Три метода}

Направление у всех одно и то же, \textbf{разные методы отличаются способом
выбора шага}:

\begin{enumerate}[label=\arabic*), leftmargin=2em]
  \item метод с дроблением шага;
  \item метод наискорейшего спуска;
  \item метод с постоянным шагом.
\end{enumerate}

Первые два — \textbf{шаговые методы}, устроенные по общей схеме:
\[
  1^\circ \ \text{находится направление спуска};
  \qquad
  2^\circ \ \text{находится шаг} .
\]

\begin{exercise}
Разобрать самостоятельно метод с постоянным шагом ($\alpha_k \equiv \alpha$).
\end{exercise}

\begin{idea}
Подсказка к упражнению: оценка~\eqref{eq:grad-main} ниже при $\alpha_k \equiv
\alpha$ сразу даёт
$f(x_{k+1}) - f(x_k) \le -\alpha(1-\alpha R)\norm{\nabla f(x_k)}^2$, так что
метод работает при $0 < \alpha < 1/R$, а наилучшая гарантия получается при
$\alpha = 1/(2R)$. Недостаток очевиден: нужна константа Липшица $R$, которую
на практике не знают.
\end{idea}

\subsubsection*{Градиентный метод с дроблением шага}

Итерация:
\[
  x_0 \ \longrightarrow\ x_{k+1} = x_k - \alpha_k \nabla f(x_k) .
\]
Шаг $\alpha_k$ подбирается \emph{дроблением}: начиная с $\alpha^{(0)}$,
пробное значение умножается на $\lambda \in (0,1)$ до тех пор, пока не будет
выполнено условие достаточного убывания.

\begin{algo}[title={Градиентный метод с дроблением шага}]
\textbf{Вход:} $x_0$, начальный шаг $\alpha^{(0)}$, коэффициент дробления
$\lambda \in (0,1)$, параметр $\eps \in (0,1)$, точность $\delta$.

\begin{steps}
  \item Если $\norm{\nabla f(x_k)} < \delta$ — остановиться.
  \item Положить $\alpha := \alpha^{(0)}$.
  \item Вычислить $x = x_k - \alpha \nabla f(x_k)$ и проверить условие
        дробления
        \[
          f(x) - f(x_k) \ \le\ -\eps\,\alpha\, \norm{\nabla f(x_k)}^2 .
        \]
  \item Если условие не выполнено — положить $\alpha := \alpha\lambda$ и
        вернуться к шагу~3.
  \item Иначе принять
        $\alpha_k = \alpha = \alpha^{(0)}\lambda^{\,i_k}$ (где $i_k$ — число
        дроблений) и $x_{k+1} = x$.
  \item Увеличить $k$ и вернуться к шагу~1.
\end{steps}
\end{algo}

\inputfig{figures/fig-step-halving}

\begin{idea}
Зачем нужен множитель $\eps$: без него условие $f(x) < f(x_k)$ выполнялось бы
при сколь угодно малом $\alpha$, и метод мог бы «застрять», делая всё более
мелкие шаги с всё меньшим выигрышем. Условие с $-\eps\alpha\norm{\nabla f}^2$
требует убывания, \emph{пропорционального} сделанному шагу, — именно это
даст в теореме оценку $\alpha_k$ снизу.
\end{idea}

\subsubsection*{Градиентный метод наискорейшего спуска}

Здесь шаг выбирается точно — решением задачи одномерной минимизации
(\S9--\S10):
\[
  \big(f(x_{k+1}) - f(x_k)\big) \longrightarrow \min_{\alpha}
  \quad\Longleftrightarrow\quad
  f(x_{k+1}) \longrightarrow \min_{\alpha}
\]
\[
  \Longleftrightarrow\quad
  f\big(x_k - \alpha \nabla f(x_k)\big) \longrightarrow \min_{\alpha \ge 0} .
\]

\begin{algo}[title={Градиентный метод наискорейшего спуска}]
\textbf{Вход:} $x_0$, точность $\delta$.

\begin{steps}
  \item Если $\norm{\nabla f(x_k)} < \delta$ — остановиться.
  \item Положить $p_k = -\nabla f(x_k)$.
  \item Решить задачу одномерной минимизации
        \[
          \alpha_k = \arg\min_{\alpha \ge 0} f(x_k + \alpha p_k)
        \]
        любым методом из \S10.
  \item Положить $x_{k+1} = x_k + \alpha_k p_k$, увеличить $k$ и вернуться
        к шагу~1.
\end{steps}
\end{algo}

\inputfig{figures/fig-zigzag}

\begin{idea}
Переход «$f(x_{k+1}) - f(x_k) \to \min$» $\iff$ «$f(x_{k+1}) \to \min$»
законен потому, что $f(x_k)$ от $\alpha$ не зависит — это просто константа,
сдвигающая минимизируемую функцию.
\end{idea}

\subsubsection*{Основная оценка}

Обе теоремы о сходимости опираются на одно неравенство, поэтому выведем его
отдельно.

\begin{lemma}\label{lem:grad-main}
Пусть $f \in C^1(\R^n)$ и градиент липшицев:
\[
  \norm{\nabla f(x) - \nabla f(y)} \le R\,\norm{x-y},
  \qquad R = \mathrm{const} > 0 .
\]
Тогда для любого $\alpha > 0$
\begin{equation}\label{eq:grad-main}
  f\big(x_k - \alpha \nabla f(x_k)\big) - f(x_k)
  \ \le\ -\alpha\,(1 - \alpha R)\, \norm{\nabla f(x_k)}^2 .
\end{equation}
\end{lemma}

\begin{proof}
Положим $x = x_k - \alpha \nabla f(x_k)$, так что
$x - x_k = -\alpha\nabla f(x_k)$. По формуле конечных приращений
\[
  f(x) - f(x_k) = \nabla\T f(x_{\mathrm{cp}})\,(x - x_k),
  \qquad
  x_{\mathrm{cp}} = x_k + \theta\,(x - x_k), \quad \theta \in [0,1] .
\]
Прибавим и вычтем $\nabla\T f(x_k)(x-x_k)$:
\[
  f(x) - f(x_k)
  = \underbrace{\nabla\T f(x_k)\,(x-x_k)}_{=\ -\alpha\norm{\nabla f(x_k)}^2}
  + \big(\nabla f(x_{\mathrm{cp}}) - \nabla f(x_k)\big)\T (x - x_k) .
\]
Второе слагаемое оцениваем неравенством Коши—Буняковского и условием
Липшица:
\[
  \big(\nabla f(x_{\mathrm{cp}}) - \nabla f(x_k)\big)\T (x - x_k)
  \le \norm{\nabla f(x_{\mathrm{cp}}) - \nabla f(x_k)}\cdot\norm{x-x_k}
  \le R\,\norm{x_{\mathrm{cp}} - x_k}\cdot\norm{x - x_k} .
\]
Но $\norm{x_{\mathrm{cp}} - x_k} = \theta\norm{x-x_k} \le \norm{x-x_k}$,
поэтому оценка продолжается до
\[
  R\,\norm{x-x_k}^2 = R\,\alpha^2 \norm{\nabla f(x_k)}^2 .
\]
Складывая,
\[
  f(x) - f(x_k)
  \le -\alpha\norm{\nabla f(x_k)}^2 + R\alpha^2\norm{\nabla f(x_k)}^2
  = \norm{\nabla f(x_k)}^2\,\big(R\alpha^2 - \alpha\big) ,
\]
что и есть~\eqref{eq:grad-main}.
\end{proof}

\subsubsection*{Сходимость метода с дроблением шага}

\begin{theorem}\label{thm:grad-split}
Пусть
\begin{itemize}[leftmargin=1.6em, topsep=2pt]
  \item $f \in C^1(\R^n)$;
  \item $f$ ограничена снизу;
  \item градиент липшицев с константой $R > 0$.
\end{itemize}
Пусть $\{x_k\}$ строится градиентным методом с дроблением шага. Тогда
\[
  \norm{\nabla f(x_k)} \xrightarrow[k \to \infty]{} 0
\]
независимо от выбора начального приближения.
\end{theorem}

\begin{proof}
\textbf{Шаг 1. Дробление конечно.} Условие приёма шага
\begin{equation}\label{eq:armijo}
  f(x) - f(x_k) \le -\eps\,\alpha\,\norm{\nabla f(x_k)}^2
\end{equation}
заведомо выполнено, если правая часть~\eqref{eq:grad-main} не превосходит
правой части~\eqref{eq:armijo}:
\[
  -\alpha(1-\alpha R) \le -\eps\alpha
  \iff 1 - \alpha R \ge \eps
  \iff \alpha \le \frac{1-\eps}{R} .
\]
Значит после конечного числа дроблений пробный шаг попадёт в этот диапазон и
будет принят.

\medskip
\textbf{Шаг 2. Монотонность.} Беря $\alpha = \alpha_k$, $x = x_{k+1}$,
получаем из~\eqref{eq:armijo}
\[
  f(x_{k+1}) - f(x_k) \le -\eps\,\alpha_k \norm{\nabla f(x_k)}^2 < 0 ,
\]
то есть $f(x_k)$ строго убывает. Ограниченная снизу убывающая
последовательность сходится, поэтому
\begin{equation}\label{eq:grad-diff0}
  f(x_k) - f(x_{k+1}) \xrightarrow[k \to \infty]{} 0 .
\end{equation}

\medskip
\textbf{Шаг 3. Оценка градиента сверху.} Переписывая~\eqref{eq:armijo}:
\begin{equation}\label{eq:grad-upper}
  \norm{\nabla f(x_k)}^2 \ \le\ \frac{f(x_k) - f(x_{k+1})}{\eps\,\alpha_k} .
\end{equation}

\medskip
\textbf{Шаг 4. Оценка шага снизу.} Пусть $i_k \ge 1$, то есть хотя бы одно
дробление было. Тогда предыдущее пробное значение
$\widetilde\alpha = \alpha_k/\lambda$ было \emph{отвергнуто}:
\[
  f\big(x_k - \widetilde\alpha \nabla f(x_k)\big) - f(x_k)
  \ >\ -\eps\,\widetilde\alpha\,\norm{\nabla f(x_k)}^2 .
\]
С другой стороны, по лемме~\ref{lem:grad-main} левая часть не превосходит
$-\widetilde\alpha(1 - \widetilde\alpha R)\norm{\nabla f(x_k)}^2$. Сравнивая,
\[
  -\widetilde\alpha(1 - \widetilde\alpha R) > -\eps\,\widetilde\alpha
  \iff 1 - \widetilde\alpha R < \eps
  \iff \widetilde\alpha > \frac{1-\eps}{R} ,
\]
откуда
\[
  \alpha_k = \lambda\,\widetilde\alpha > \lambda\,\frac{1-\eps}{R} .
\]
Если же $i_k = 0$, то $\alpha_k = \alpha^{(0)}$. В обоих случаях
\begin{equation}\label{eq:alpha-lower}
  \alpha_k \ \ge\ \bar\alpha := \min\left\{\alpha^{(0)},\ \lambda\,\frac{1-\eps}{R}\right\} > 0 .
\end{equation}

\medskip
\textbf{Шаг 5. Завершение.} Подставляя~\eqref{eq:alpha-lower}
в~\eqref{eq:grad-upper}:
\[
  \norm{\nabla f(x_k)}^2
  \ \le\ \frac{f(x_k) - f(x_{k+1})}{\eps\,\bar\alpha}
  \xrightarrow[k \to \infty]{} 0
\]
в силу~\eqref{eq:grad-diff0}. Значит $\norm{\nabla f(x_k)} \to 0$. Начальное
приближение нигде не использовалось, то есть сходимость глобальная.
\end{proof}

\begin{idea}[title={Зачем нужны обе оценки шага}]
В тетради доказательство обрывается на словах «\emph{одновременно получили
верхнюю и нижнюю оценки}» со знаками вопроса; шаг~5 восстановлен.

\medskip
Смысл восклицания такой: оценка $\alpha \le \frac{1-\eps}{R}$
\emph{сверху} гарантирует, что дробление когда-нибудь закончится, а оценка
$\alpha_k > \lambda\frac{1-\eps}{R}$ \emph{снизу} гарантирует, что принятый
шаг не окажется слишком мелким. Вторая нужна именно для того, чтобы в
неравенстве~\eqref{eq:grad-upper} можно было избавиться от $\alpha_k$ в
знаменателе: без неё дробь могла бы не стремиться к нулю.

\medskip
Случай $i_k = 0$ в тетради не оговорён — там выкладка неявно считает, что
хотя бы одно дробление произошло. Это дополнено выше.
\end{idea}

\begin{idea}
Обратите внимание, что доказано: стремление к нулю \emph{градиента}, а не
сходимость $x_k$ к точке минимума. Без дополнительных предположений большего
и не получить — $f$ может вообще не иметь минимума (например, монотонно
убывать к своей нижней грани). Сходимость самих точек даёт следующий
параграф, где добавляется сильная выпуклость.
\end{idea}

\subsubsection*{Сходимость метода наискорейшего спуска}

\begin{theorem}\label{thm:steepest}
Пусть выполнены условия теоремы~\ref{thm:grad-split}, а $\alpha_k$ выбирается
по принципу наискорейшего спуска. Тогда
\[
  \norm{\nabla f(x_k)} \xrightarrow[k \to \infty]{} 0
\]
независимо от выбора начального приближения.
\end{theorem}

\begin{proof}
По лемме~\ref{lem:grad-main} для любого $\alpha > 0$
\[
  f\big(x_k - \alpha\nabla f(x_k)\big) - f(x_k)
  \le -\alpha(1-\alpha R)\norm{\nabla f(x_k)}^2 .
\]
Правая часть как функция $\alpha$ — парабола $R\alpha^2 - \alpha$, умноженная
на $\norm{\nabla f(x_k)}^2$; её минимум достигается при
\[
  \alpha = \frac{1}{2R},
  \qquad\text{и равен}\qquad
  -\frac{1}{4R}\,\norm{\nabla f(x_k)}^2 .
\]
Шаг $\alpha_k$ выбран из условия наименьшего значения $f$, поэтому для любого
$\alpha$
\[
  f(x_k + \alpha_k p_k) - f(x_k) \ \le\ f(x_k + \alpha p_k) - f(x_k),
  \qquad p_k = -\nabla f(x_k) .
\]
Берём $\alpha = \tfrac{1}{2R}$:
\[
  f(x_{k+1}) - f(x_k) \ \le\ -\frac{1}{4R}\,\norm{\nabla f(x_k)}^2 .
\]
Отсюда $f(x_k)$ строго убывает и ограничена снизу, значит
$f(x_k) - f(x_{k+1}) \to 0$, а
\begin{equation}\label{eq:steepest-bound}
  \norm{\nabla f(x_k)}^2 \ \le\ 4R\,\big(f(x_k) - f(x_{k+1})\big)
  \xrightarrow[k\to\infty]{} 0 .
\end{equation}
\end{proof}

\begin{idea}
Фокус доказательства в том, что \emph{знать} оптимальный шаг не требуется:
$\alpha_k$ не вычисляется явно, он лишь не хуже «пробного» $\alpha = 1/(2R)$.
Именно поэтому константа Липшица входит в оценку, но не в сам метод.
\end{idea}
